\documentclass[11pt]{article}

% ============================================================
% Encoding and page layout
% ============================================================

\usepackage[T1]{fontenc}
\usepackage[utf8]{inputenc}
\usepackage[a4paper,margin=1.12in]{geometry}

% ============================================================
% Mathematics and theorem environments
% ============================================================

\usepackage{amsmath,amssymb,amsthm,mathtools}
\usepackage{microtype}
\usepackage{xcolor}
\usepackage{enumitem}
\usepackage{hyperref}
\usepackage[nameinlink,capitalize]{cleveref}

\hypersetup{
    colorlinks=true,
    linkcolor=blue!55!black,
    citecolor=blue!55!black,
    urlcolor=blue!55!black
}

\newtheorem{theorem}{Theorem}[section]
\newtheorem{proposition}[theorem]{Proposition}
\newtheorem{lemma}[theorem]{Lemma}
\newtheorem{corollary}[theorem]{Corollary}
\newtheorem{definition}[theorem]{Definition}
\newtheorem{remark}[theorem]{Remark}
\newtheorem{assumption}[theorem]{Assumption}

% ============================================================
% Macros
% ============================================================

\newcommand{\ii}{\mathrm{i}}
\newcommand{\dd}{\,\mathrm{d}}

\newcommand{\zetaR}{\zeta}
\newcommand{\half}{\frac12}

\newcommand{\NQ}{\mathcal N_Q}
\newcommand{\RQ}{\mathcal R_Q}
\newcommand{\Rtot}{\mathcal R_{\rm total}}
\newcommand{\Beff}{\mathfrak B_{\rm eff}}

\newcommand{\Ech}{\mathcal E}
\newcommand{\Och}{\mathcal O}
\newcommand{\Dvec}{\mathfrak D}
\newcommand{\Bvec}{\mathfrak B}
\newcommand{\Rvec}{\mathfrak R}

\newcommand{\gev}{g_0}
\newcommand{\Cinner}{C_{\rm inner}}
\newcommand{\Czero}{C_0}
\newcommand{\Cone}{C_1}

\newcommand{\RepoURL}{\url{https://github.com/vicblab/RH-outer.strip-certifier}}

% ============================================================
% Title
% ============================================================

\title{
\textbf{A Reflected Two-Channel Quotient Obstruction}\\[0.35em]
\large A finite derivative-certificate reduction for inner-strip zeta zeros
}

\author{V\'ictor Blanco Bataller}

\date{\today}

\begin{document}

\maketitle

% ============================================================
% Abstract
% ============================================================

\begin{abstract}
This paper develops the inner-strip continuation of the Mellin-symmetric
quotient-obstruction method.  The outer-strip argument proves a high-height
zero-free region by comparing a quotient numerator
\[
        \NQ(s)
\]
with an effective bracket
\[
        \frac{\ii}{\tau}\Beff(\sigma,\tau).
\]
Away from the critical line, the bracket has a fixed positive lower bound.
Near the critical line, however, this one-channel lower bound degenerates.
Thus the outer-strip proof cannot be pushed to \(\Re s=1/2\) by merely
tightening constants.

The inner-strip mechanism is different.  If
\[
        s_-=\frac12-u+\ii\tau
\]
is an off-critical zero, then the functional equation and conjugation symmetry
force the reflected point
\[
        s_+=\frac12+u+\ii\tau
\]
to be a zero as well.  Hence an off-critical reflected zero pair forces two
quotient equations:
\[
        \NQ(s_-)=0,
        \qquad
        \NQ(s_+)=0.
\]
Equivalently, it forces both the reflected average and the reflected finite
difference to vanish.  This leads to the two-channel detector
\[
        \Dvec(u,\tau)
        =
        \left(
        \frac{\NQ(s_-)+\NQ(s_+)}2,
        \frac{\NQ(s_-)-\NQ(s_+)}{2u}
        \right).
\]

The main structural observation is that the corresponding bracket vector does
not degenerate at the critical line.  Its limiting components are
\[
        \Beff\left(\frac12,\tau\right)
        \quad\text{and}\quad
        -\partial_\sigma\Beff\left(\frac12,\tau\right),
\]
and these two quantities cannot vanish simultaneously.  In fact, the critical
line bracket vector has the explicit lower scale
\[
        \frac{\sqrt2}{48}(23\log2-2)>0.
\]
For a finite inner strip
\[
        0<\left|\Re s-\frac12\right|\le0.01,
\]
this lower bound is replaced by a finite interval certificate, for which the
conservative public target is
\[
        \gev=0.2.
\]

The residual side requires one new ingredient beyond the outer-strip
certificate: each residual sector must be certified not only in value but also
in its \(\sigma\)-derivative.  The ordinary residual controls the reflected
average, while the derivative residual controls the divided reflected
difference by the mean value theorem.  Thus the inner-strip problem is reduced
to a finite certificate for the seven residual sectors and their
\(\sigma\)-derivatives.

This paper is a structural reduction.  It does not by itself claim a completed
inner-strip zero-free theorem.  The certified theorem is obtained only after
the public file
\[
        \texttt{INNER\_STRIP\_CERTIFICATE\_SAFE.json}
\]
has status
\[
        \texttt{inner\_full\_certificate}.
\]
\end{abstract}

\tableofcontents

% ============================================================
% 1. Introduction
% ============================================================

\section{Introduction}
\label{sec:introduction}

The goal of this paper is to isolate the correct inner-strip obstruction for
possible off-critical zeros of the Riemann zeta function.  It is a continuation
of the outer-strip quotient-obstruction framework developed in the structural
outer-strip paper and certified in the companion finite certificate.  The
outer-strip mechanism is based on an asymptotic identity of the form
\[
        \NQ(s)
        =
        \frac{\ii}{\tau}\Beff(\sigma,\tau)
        +
        \Rtot(s),
        \qquad
        s=\sigma+\ii\tau,
\]
where \(\NQ(s)\) is a quotient-normalized numerator, \(\Beff(\sigma,\tau)\) is
an effective bracket, and \(\Rtot(s)\) is the grouped residual.

In a fixed outer strip, for example
\[
        0.01\le\sigma\le0.49,
\]
the bracket satisfies a positive lower bound
\[
        |\Beff(\sigma,\tau)|\ge d>0.
\]
The residual is certified to satisfy
\[
        |\Rtot(s)|
        \le
        C
        \frac{\log\tau}{\tau^{3/2}}.
\]
Thus the main term has size at least \(d/\tau\), while the residual has size
\(O(\log\tau/\tau^{3/2})\).  For sufficiently large \(\tau\), the main term
dominates.  Since every zero satisfies
\[
        \NQ(s)=0,
\]
this gives a contradiction and therefore excludes zeros in the fixed outer
strip.

This argument is deliberately an outer-strip argument.  It cannot be extended
to the critical line by merely improving the constant \(C\).  The reason is
structural: the one-point bracket lower bound degenerates as
\[
        \sigma\to\frac12.
\]
If
\[
        u=\frac12-\sigma,
\]
then the one-channel bracket gap is of first order in \(u\).  Consequently the
outer-strip comparison can at best see a main term of size comparable to
\[
        \frac{u}{\tau}.
\]
A residual estimate of size
\[
        \frac{\log\tau}{\tau^{3/2}}
\]
does not by itself exclude arbitrarily small positive \(u\).  Therefore the
inner strip requires a different detector.

The correct detector comes from the symmetry of the zeros.  Suppose
\[
        s_-=\frac12-u+\ii\tau,
        \qquad
        0<u\le u_0,
\]
is a nontrivial zero.  By the functional equation and conjugation symmetry,
the reflected point
\[
        s_+=\frac12+u+\ii\tau
\]
is also a nontrivial zero.  Hence the zero condition gives two exact equations:
\[
        \NQ(s_-)=0,
        \qquad
        \NQ(s_+)=0.
\]
It is therefore wasteful to test only one value of \(\NQ\).  The reflected pair
forces both the average
\[
        \frac{\NQ(s_-)+\NQ(s_+)}2
\]
and the divided difference
\[
        \frac{\NQ(s_-)-\NQ(s_+)}{2u}
\]
to vanish.

This is the key inner-strip idea.  The outer-strip bracket weakens at the
critical line, but an off-critical reflected pair is more rigid than a single
point: it must make both the reflected midpoint value and the reflected slope
vanish.  The two-channel detector records precisely these two quantities.

We therefore define
\[
        \Dvec(u,\tau)
        :=
        \left(
        \Ech(u,\tau),
        \Och(u,\tau)
        \right),
\]
where
\[
        \Ech(u,\tau)
        :=
        \frac{\NQ(s_-)+\NQ(s_+)}2,
\]
and
\[
        \Och(u,\tau)
        :=
        \frac{\NQ(s_-)-\NQ(s_+)}{2u}.
\]
At an off-critical reflected zero pair,
\[
        \Dvec(u,\tau)=(0,0).
\]

Substituting the structural quotient expansion into the two channels gives
\[
        \Dvec(u,\tau)
        =
        \frac{\ii}{\tau}\Bvec(u,\tau)
        +
        \Rvec(u,\tau),
\]
where \(\Bvec\) is the reflected bracket vector and \(\Rvec\) is the reflected
residual vector.  The bracket vector is
\[
        \Bvec(u,\tau)
        =
        \left(
        B_{\rm ev}(u,\tau),
        B_{\rm od}(u,\tau)
        \right),
\]
with
\[
        B_{\rm ev}(u,\tau)
        =
        \frac{
        \Beff(\frac12-u,\tau)
        +
        \Beff(\frac12+u,\tau)
        }2,
\]
and
\[
        B_{\rm od}(u,\tau)
        =
        \frac{
        \Beff(\frac12-u,\tau)
        -
        \Beff(\frac12+u,\tau)
        }{2u}.
\]
As \(u\to0\), these components tend to
\[
        \Beff\left(\frac12,\tau\right)
        \quad\text{and}\quad
        -\partial_\sigma\Beff\left(\frac12,\tau\right).
\]
The crucial point is that these limiting components cannot vanish
simultaneously.  The reflected even channel sees the value of the bracket on
the critical line; the reflected odd channel sees its transverse derivative.

This is the inner-strip replacement for the outer-strip bracket gap.  Instead
of proving
\[
        |\Beff(\sigma,\tau)|\ge d
\]
at one point, we prove a vector lower bound
\[
        \|\Bvec(u,\tau)\|\ge \gev.
\]
For the final certified inner strip
\[
        0<u\le0.01,
\]
the conservative target is
\[
        \gev=0.2.
\]

The residual side also changes.  The even residual is controlled by the
ordinary residual bound.  The divided odd residual is controlled by a
\(\sigma\)-derivative bound.  Thus each residual sector must be certified in
two norms:
\[
        |\mathcal R_\nu(s)|
        \le
        C_{\nu,0}
        \frac{\log\tau}{\tau^{3/2}},
\]
and
\[
        |\partial_\sigma\mathcal R_\nu(s)|
        \le
        C_{\nu,1}
        \frac{\log\tau}{\tau^{3/2}}.
\]
The seven sectors are the same sectors as in the outer-strip decomposition:
\[
        {\rm band},
        \quad
        {\rm bd},
        \quad
        {\rm far},
        \quad
        {\rm end},
        \quad
        {\rm nonstat},
        \quad
        {\rm floor},
        \quad
        {\rm core}.
\]
The derivative certificate is what makes the divided reflected difference
finite and uniform as \(u\to0\).

The resulting exclusion mechanism is simple.  If
\[
        \|\Bvec(u,\tau)\|\ge \gev
\]
and
\[
        \|\Rvec(u,\tau)\|
        \le
        \Cinner
        \frac{\log\tau}{\tau^{3/2}},
\]
then
\[
        \|\Dvec(u,\tau)\|
        \ge
        \frac{\gev}{\tau}
        -
        \Cinner
        \frac{\log\tau}{\tau^{3/2}}.
\]
This is positive whenever
\[
        \sqrt{\tau}
        >
        \frac{\Cinner}{\gev}
        \log\tau.
\]
But a reflected off-critical zero pair would force
\[
        \Dvec(u,\tau)=(0,0).
\]
Hence no such pair can exist above the certified height.

The paper is organized as follows.  In
\cref{sec:structural-input}, we state the structural quotient expansion
imported from the outer-strip paper.  In
\cref{sec:reflected-zero-pairs}, we record the exact reflected zero-pair
identity.  In \cref{sec:two-channel-detector}, we define the two-channel
detector and derive its bracket-residual decomposition.  In
\cref{sec:bracket-vector}, we prove the algebraic nondegeneracy of the
critical-line bracket vector and state the finite inner-strip bracket
certificate.  In \cref{sec:residual-vector}, we show how ordinary and
\(\sigma\)-derivative residual certificates imply the two-channel residual
bound.  In \cref{sec:sector-certifiability}, we explain sector by sector why
the derivative constants are finitely certifiable.  Finally, in
\cref{sec:inner-exclusion}, we state the certified inner-strip exclusion
theorem conditional on the final public certificate
\[
        \texttt{INNER\_STRIP\_CERTIFICATE\_SAFE.json}.
\]

\begin{remark}[Scope]
This paper is an inner-strip structural reduction.  It does not claim a
completed inner-strip zero-free theorem until the derivative-aware certificate
has been produced and the final JSON file has status
\[
        \texttt{inner\_full\_certificate}.
\]
The purpose of the paper is to prove that the inner-strip problem reduces to a
finite list of ordinary and \(\sigma\)-derivative residual certificates.
\end{remark}

% ============================================================
% 2. Structural input from the outer-strip quotient framework
% ============================================================

\section{Structural input from the outer-strip quotient framework}
\label{sec:structural-input}

This section records the analytic input imported from the outer-strip
quotient-obstruction framework.  The inner-strip argument does not rederive the
Mellin identities, the four-family decomposition, the Poisson band expansion,
or the compact-boundary construction.  Those belong to the structural
outer-strip paper.  What changes here is the way in which the already-derived
quotient expansion is tested near the critical line.

Throughout this paper we write
\[
        s=\sigma+\ii\tau,
        \qquad
        \tau>0,
\]
and we use the quotient numerator
\[
        \NQ(s)
\]
from the outer-strip paper.  At every nontrivial zero of \(\zeta(s)\),
\[
        \NQ(s)=0.
\]
The precise normalization used here is the same as in the outer-strip paper:
\[
        \NQ(s)
        =
        B(s)+\frac{s-1}{s}A(s),
\]
where
\[
        A(s):=\int_1^\infty \{x\}x^{-s-1}\dd x,
\]
and
\[
        B(s):=\int_1^\infty \{x\}x^{s-2}\dd x.
\]
Equivalently, one may use the multiplied numerator
\[
        sB(s)+(s-1)A(s),
\]
since \(s\ne0\) in the critical strip.  In this paper we use the quotient
normalization \(\NQ\), because it is the normalization in which the effective
bracket appears with prefactor \(\ii/\tau\).

\subsection{Imported quotient expansion}

The structural outer-strip paper proves an expansion of the form
\[
        \NQ(s)
        =
        \frac{\ii}{\tau}\Beff(\sigma,\tau)
        +
        \Rtot(s),
\]
where \(\Beff(\sigma,\tau)\) is the effective quotient bracket and
\(\Rtot(s)\) is the total grouped residual.

For the inner-strip argument we require this expansion on the symmetric
neighborhood
\[
        \left|\sigma-\frac12\right|\le0.01,
        \qquad
        \tau\ge T_0,
\]
where in the certified implementation one may take
\[
        T_0=10.
\]
The logarithm in all residual norms is the natural logarithm.

\begin{assumption}[Imported quotient expansion on the inner strip]
\label{ass:inner-quotient-expansion}
On the region
\[
        \left|\sigma-\frac12\right|\le0.01,
        \qquad
        \tau\ge T_0,
\]
the quotient numerator admits the decomposition
\[
        \NQ(s)
        =
        \frac{\ii}{\tau}\Beff(\sigma,\tau)
        +
        \Rtot(s).
\]
The total residual decomposes into the seven sectors
\[
        \Rtot(s)
        =
        \mathcal R_{\rm band}(s)
        +
        \mathcal R_{\rm bd}(s)
        +
        \mathcal R_{\rm far}(s)
        +
        \mathcal R_{\rm end}(s)
        +
        \mathcal R_{\rm nonstat}(s)
        +
        \mathcal R_{\rm floor}(s)
        +
        \mathcal R_{\rm core}(s).
\]
The sector names and normalizations are the same as in the outer-strip
structural paper and its certification companion.
\end{assumption}

\begin{remark}[Why this is an assumption here]
The present paper is the inner-strip structural reduction.  It is not intended
to repeat the analytic derivation of the quotient expansion.  The role of
\cref{ass:inner-quotient-expansion} is to state the exact interface required
from the outer-strip structural paper.  The new work here begins after this
interface: we replace the one-channel outer-strip test by a reflected
two-channel test and identify the additional derivative residual certificates
needed near the critical line.
\end{remark}

\subsection{Effective bracket}

The effective bracket has the form
\[
\begin{aligned}
        \Beff(\sigma,\tau)
        &=
        2^\sigma e^{\ii L}
        \left(
        \frac{12-\sigma}{48}
        +
        \frac{\ii X}{4}
        \right)
\\
        &\quad
        -
        2^{-\sigma}e^{-\ii L}
        \left(
        \frac{11+\sigma}{24}
        -
        \frac{\ii X}{2}
        \right),
\end{aligned}
\]
where
\[
        L:=\tau\log2,
\]
and
\[
        X:=\Xi_{\rm end}(\tau)
\]
is the real regularized endpoint scalar imported from the structural
four-family endpoint analysis.

The only properties of \(X\) used in the bracket-vector argument are that
\[
        X\in\mathbb R
\]
and that the same value of \(X\) appears in the reflected evaluations
\[
        \Beff\left(\frac12-u,\tau\right),
        \qquad
        \Beff\left(\frac12+u,\tau\right).
\]
The critical bracket-vector lower bound proved later is uniform in real \(X\);
therefore no numerical upper bound for \(X\) is needed in the algebraic
critical-line nondegeneracy argument.

\begin{remark}[Why no bound for \(X\) is needed at the critical line]
The outer-strip proof uses a one-channel modulus lower bound for
\(\Beff(\sigma,\tau)\).  That lower bound depends on comparing two amplitudes
and degenerates as \(\sigma\to1/2\).  The inner-strip proof instead uses the
two quantities
\[
        \Beff\left(\frac12,\tau\right)
        \quad\text{and}\quad
        \partial_\sigma\Beff\left(\frac12,\tau\right).
\]
Their combined squared norm has a positive lower bound uniformly for every
real value of \(X\).  This is the algebraic reason the reflected two-channel
detector is the correct inner-strip object.
\end{remark}

\subsection{Residual norms required for the inner strip}

The outer-strip certificate controls the residual itself in the norm
\[
        \frac{\log\tau}{\tau^{3/2}}.
\]
For the inner strip, this is not enough.  The reflected average channel uses
the ordinary residual, but the divided odd channel contains
\[
        \frac{\Rtot(s_-)-\Rtot(s_+)}{2u}.
\]
To control this uniformly as \(u\to0\), one needs a uniform
\(\sigma\)-derivative bound for the residual.

Therefore each residual sector must be certified in two norms.

\begin{definition}[Inner-strip value and derivative sector constants]
\label{def:inner-sector-constants}
Let
\[
        \nu\in
        \{
        {\rm band},
        {\rm bd},
        {\rm far},
        {\rm end},
        {\rm nonstat},
        {\rm floor},
        {\rm core}
        \}.
\]
A pair of constants
\[
        C_{\nu,0},
        \qquad
        C_{\nu,1}
\]
certifies the sector \(\mathcal R_\nu\) on the inner strip if, for every
\[
        \left|\sigma-\frac12\right|\le0.01,
        \qquad
        \tau\ge T_0,
\]
one has
\[
        |\mathcal R_\nu(\sigma+\ii\tau)|
        \le
        C_{\nu,0}
        \frac{\log\tau}{\tau^{3/2}},
\]
and
\[
        |\partial_\sigma\mathcal R_\nu(\sigma+\ii\tau)|
        \le
        C_{\nu,1}
        \frac{\log\tau}{\tau^{3/2}}.
\]
\end{definition}

Define the total value and derivative constants by
\[
        \Czero
        :=
        \sum_\nu C_{\nu,0},
\]
and
\[
        \Cone
        :=
        \sum_\nu C_{\nu,1}.
\]
Finally define the inner vector residual constant
\[
        \Cinner
        :=
        \sqrt{\Czero^2+\Cone^2}.
\]

\begin{remark}[Why the derivative norm is the right one]
The divided odd channel contains a difference quotient in the real direction.
For
\[
        s_\pm=\frac12\pm u+\ii\tau,
\]
the mean value theorem gives
\[
        \left|
        \frac{\Rtot(s_-)-\Rtot(s_+)}{2u}
        \right|
        \le
        \sup_{\left|\sigma-\frac12\right|\le u}
        |\partial_\sigma\Rtot(\sigma+\ii\tau)|.
\]
Thus a uniform \(\partial_\sigma\)-residual certificate is exactly what is
needed to make the divided odd channel finite and stable as \(u\to0\).
\end{remark}

\subsection{Certification interface}

The final inner-strip theorem will not cite individual hand-transcribed
constants.  Instead, it will cite a public JSON certificate.  The required
certificate contract is as follows.

For each sector \(\nu\), the corresponding JSON file must have status
\[
        \texttt{proved}
\]
and must contain the two fields
\[
        \texttt{C\_value\_upper\_safe},
        \qquad
        \texttt{C\_sigma\_derivative\_upper\_safe}.
\]
The final merge file
\[
        \texttt{INNER\_STRIP\_CERTIFICATE\_SAFE.json}
\]
must have status
\[
        \texttt{inner\_full\_certificate}.
\]
It must record
\[
        \gev,
        \qquad
        \Czero,
        \qquad
        \Cone,
        \qquad
        \Cinner,
\]
and the verified height inequality
\[
        \sqrt T
        >
        \frac{\Cinner}{\gev}\log T.
\]

For the public target used in this paper, the bracket-vector certificate will
use
\[
        \gev=0.2.
\]
If the final merged certificate verifies
\[
        T=10^{30},
\]
then the inner-strip exclusion height matches the outer-strip certified height.

\begin{remark}[No hidden constants]
All constants used in the final inner-strip theorem must appear either in the
structural formulas of this paper or in the final public certificate file.
There is no untracked residual sector and no untracked derivative sector.
\end{remark}

% ============================================================
% 3. Reflected zero pairs and the two-channel detector
% ============================================================

\section{Reflected zero pairs}
\label{sec:reflected-zero-pairs}

The inner-strip argument begins with an exact symmetry observation.  In the
outer strip, a single point \(s=\sigma+\ii\tau\) is tested against the
one-point bracket \(\Beff(\sigma,\tau)\).  Near the critical line this is not
the correct perspective.  A zero away from the critical line never occurs in
isolation: by the functional equation and conjugation symmetry it belongs to a
reflected pair.

We write
\[
        s_-(u,\tau)
        :=
        \frac12-u+\ii\tau,
        \qquad
        s_+(u,\tau)
        :=
        \frac12+u+\ii\tau,
\]
where
\[
        0<u\le u_0.
\]
In the certified inner strip we shall take
\[
        u_0=0.01.
\]

\begin{lemma}[Reflected zero-pair symmetry]
\label{lem:reflected-zero-pair}
Let
\[
        0<u<\frac12,
        \qquad
        \tau>0.
\]
If
\[
        s_-=\frac12-u+\ii\tau
\]
is a nontrivial zero of \(\zeta(s)\), then
\[
        s_+=\frac12+u+\ii\tau
\]
is also a nontrivial zero of \(\zeta(s)\).
\end{lemma}

\begin{proof}
The nontrivial zeros of \(\zeta(s)\) are invariant under conjugation
\[
        s\mapsto \overline{s}
\]
and under the functional-equation reflection
\[
        s\mapsto 1-s.
\]
Equivalently, they are invariant under the combined reflection
\[
        s\mapsto 1-\overline{s}.
\]
Applying this map to
\[
        s_-=\frac12-u+\ii\tau
\]
gives
\[
        1-\overline{s_-}
        =
        1-\left(\frac12-u-\ii\tau\right)
        =
        \frac12+u+\ii\tau
        =
        s_+.
\]
Therefore \(s_+\) is also a nontrivial zero.
\end{proof}

The quotient numerator was normalized in the structural paper so that every
nontrivial zero satisfies
\[
        \NQ(s)=0.
\]
The preceding symmetry therefore gives two exact equations from one
off-critical zero.

\begin{corollary}[Two quotient equations at a reflected pair]
\label{cor:two-quotient-equations}
If
\[
        s_-=\frac12-u+\ii\tau
\]
is a nontrivial zero with \(0<u<1/2\), then
\[
        \NQ(s_-)=0,
        \qquad
        \NQ(s_+)=0,
\]
where
\[
        s_+=\frac12+u+\ii\tau.
\]
\end{corollary}

\begin{proof}
By \cref{lem:reflected-zero-pair}, \(s_+\) is also a nontrivial zero.  The
quotient numerator satisfies \(\NQ(s)=0\) at every nontrivial zero.  Applying
this to both \(s_-\) and \(s_+\) proves the claim.
\end{proof}

\begin{remark}[Why this matters]
The outer-strip proof uses the single equation
\[
        \NQ(s)=0
\]
at one point.  In the inner strip, an off-critical zero gives the stronger
pair of equations
\[
        \NQ(s_-)=0,
        \qquad
        \NQ(s_+)=0.
\]
The two-channel detector below is simply the average and the finite-difference
form of these two equations.
\end{remark}

% ============================================================
% 4. The two-channel detector
% ============================================================

\section{The two-channel detector}
\label{sec:two-channel-detector}

We now define the reflected detector used throughout the inner-strip argument.
It is useful to think of the two components as follows:

\begin{itemize}[leftmargin=2.2em]
    \item the even channel measures the reflected midpoint value of the
    quotient numerator;

    \item the divided odd channel measures the reflected transverse slope of
    the quotient numerator.
\end{itemize}

At an off-critical reflected zero pair, both quantities must vanish.

\begin{definition}[Even and divided odd quotient channels]
\label{def:two-channel-detector}
Let
\[
        s_-=\frac12-u+\ii\tau,
        \qquad
        s_+=\frac12+u+\ii\tau,
\]
with
\[
        0<u\le u_0.
\]
Define the even quotient channel by
\[
        \Ech(u,\tau)
        :=
        \frac{\NQ(s_-)+\NQ(s_+)}{2},
\]
and the divided odd quotient channel by
\[
        \Och(u,\tau)
        :=
        \frac{\NQ(s_-)-\NQ(s_+)}{2u}.
\]
The reflected two-channel detector is
\[
        \Dvec(u,\tau)
        :=
        \left(
        \Ech(u,\tau),
        \Och(u,\tau)
        \right).
\]
\end{definition}

\begin{lemma}[Detector vanishing at an off-critical reflected zero pair]
\label{lem:detector-vanishing}
If \(s_-=\frac12-u+\ii\tau\) is a nontrivial zero with \(0<u<1/2\), then
\[
        \Dvec(u,\tau)=(0,0).
\]
\end{lemma}

\begin{proof}
By \cref{cor:two-quotient-equations},
\[
        \NQ(s_-)=0,
        \qquad
        \NQ(s_+)=0.
\]
Therefore
\[
        \Ech(u,\tau)
        =
        \frac{0+0}{2}
        =
        0,
\]
and
\[
        \Och(u,\tau)
        =
        \frac{0-0}{2u}
        =
        0.
\]
Thus
\[
        \Dvec(u,\tau)=(0,0).
\]
\end{proof}

\subsection{Bracket and residual channels}

We now insert the structural expansion from
\cref{ass:inner-quotient-expansion} into the two reflected channels.  For
\[
        s_\pm=\frac12\pm u+\ii\tau,
\]
the expansion gives
\[
        \NQ(s_\pm)
        =
        \frac{\ii}{\tau}
        \Beff\left(\frac12\pm u,\tau\right)
        +
        \Rtot(s_\pm).
\]

Define the even bracket channel by
\[
        B_{\rm ev}(u,\tau)
        :=
        \frac{
        \Beff(\frac12-u,\tau)
        +
        \Beff(\frac12+u,\tau)
        }{2},
\]
and the divided odd bracket channel by
\[
        B_{\rm od}(u,\tau)
        :=
        \frac{
        \Beff(\frac12-u,\tau)
        -
        \Beff(\frac12+u,\tau)
        }{2u}.
\]
The bracket vector is
\[
        \Bvec(u,\tau)
        :=
        \left(
        B_{\rm ev}(u,\tau),
        B_{\rm od}(u,\tau)
        \right).
\]

Similarly, define the even residual channel by
\[
        R_{\rm ev}(u,\tau)
        :=
        \frac{
        \Rtot(s_-)
        +
        \Rtot(s_+)
        }{2},
\]
and the divided odd residual channel by
\[
        R_{\rm od}(u,\tau)
        :=
        \frac{
        \Rtot(s_-)
        -
        \Rtot(s_+)
        }{2u}.
\]
The residual vector is
\[
        \Rvec(u,\tau)
        :=
        \left(
        R_{\rm ev}(u,\tau),
        R_{\rm od}(u,\tau)
        \right).
\]

\begin{proposition}[Two-channel bracket-residual decomposition]
\label{prop:two-channel-decomposition}
Assume \cref{ass:inner-quotient-expansion}.  For
\[
        0<u\le u_0,
        \qquad
        \tau\ge T_0,
\]
one has
\[
        \Dvec(u,\tau)
        =
        \frac{\ii}{\tau}
        \Bvec(u,\tau)
        +
        \Rvec(u,\tau).
\]
Equivalently,
\[
        \Ech(u,\tau)
        =
        \frac{\ii}{\tau}
        B_{\rm ev}(u,\tau)
        +
        R_{\rm ev}(u,\tau),
\]
and
\[
        \Och(u,\tau)
        =
        \frac{\ii}{\tau}
        B_{\rm od}(u,\tau)
        +
        R_{\rm od}(u,\tau).
\]
\end{proposition}

\begin{proof}
Using \cref{ass:inner-quotient-expansion} at \(s_-\) and \(s_+\), we obtain
\[
        \NQ(s_-)
        =
        \frac{\ii}{\tau}
        \Beff\left(\frac12-u,\tau\right)
        +
        \Rtot(s_-),
\]
and
\[
        \NQ(s_+)
        =
        \frac{\ii}{\tau}
        \Beff\left(\frac12+u,\tau\right)
        +
        \Rtot(s_+).
\]
Adding these two identities and dividing by \(2\) gives
\[
        \Ech(u,\tau)
        =
        \frac{\ii}{\tau}
        \frac{
        \Beff(\frac12-u,\tau)
        +
        \Beff(\frac12+u,\tau)
        }2
        +
        \frac{\Rtot(s_-)+\Rtot(s_+)}2.
\]
By definition, this is
\[
        \Ech(u,\tau)
        =
        \frac{\ii}{\tau}
        B_{\rm ev}(u,\tau)
        +
        R_{\rm ev}(u,\tau).
\]

Subtracting the two quotient expansions and dividing by \(2u\) gives
\[
        \Och(u,\tau)
        =
        \frac{\ii}{\tau}
        \frac{
        \Beff(\frac12-u,\tau)
        -
        \Beff(\frac12+u,\tau)
        }{2u}
        +
        \frac{
        \Rtot(s_-)
        -
        \Rtot(s_+)
        }{2u}.
\]
By definition, this is
\[
        \Och(u,\tau)
        =
        \frac{\ii}{\tau}
        B_{\rm od}(u,\tau)
        +
        R_{\rm od}(u,\tau).
\]
Combining the two component identities proves the vector identity.
\end{proof}

\begin{remark}[The divided odd channel is harmless only after derivative certification]
The expression
\[
        \frac{\Rtot(s_-)-\Rtot(s_+)}{2u}
\]
is not controlled uniformly by an ordinary residual bound alone.  Uniformity
as \(u\to0\) requires a bound for
\[
        \partial_\sigma \Rtot(\sigma+\ii\tau).
\]
This is why the inner-strip certificate must include \(\sigma\)-derivative
constants for every residual sector.
\end{remark}

\subsection{The contradiction mechanism}

The proof of the inner-strip exclusion will have the following form.  Suppose
that an off-critical reflected zero pair exists.  By
\cref{lem:detector-vanishing},
\[
        \Dvec(u,\tau)=(0,0).
\]
On the other hand, by \cref{prop:two-channel-decomposition},
\[
        \Dvec(u,\tau)
        =
        \frac{\ii}{\tau}\Bvec(u,\tau)
        +
        \Rvec(u,\tau).
\]
Therefore
\[
        \|\Dvec(u,\tau)\|
        \ge
        \frac{\|\Bvec(u,\tau)\|}{\tau}
        -
        \|\Rvec(u,\tau)\|.
\]
If we prove
\[
        \|\Bvec(u,\tau)\|\ge \gev
\]
and
\[
        \|\Rvec(u,\tau)\|
        \le
        \Cinner
        \frac{\log\tau}{\tau^{3/2}},
\]
then
\[
        \|\Dvec(u,\tau)\|
        \ge
        \frac{\gev}{\tau}
        -
        \Cinner
        \frac{\log\tau}{\tau^{3/2}}.
\]
The right-hand side is positive whenever
\[
        \sqrt{\tau}
        >
        \frac{\Cinner}{\gev}
        \log\tau.
\]
This contradicts
\[
        \Dvec(u,\tau)=(0,0).
\]

Thus the entire inner-strip problem reduces to two tasks:

\begin{enumerate}[label=(\roman*)]
    \item prove or certify the bracket-vector lower bound
    \[
        \|\Bvec(u,\tau)\|\ge \gev;
    \]

    \item prove or certify the residual-vector upper bound
    \[
        \|\Rvec(u,\tau)\|
        \le
        \Cinner
        \frac{\log\tau}{\tau^{3/2}}.
    \]
\end{enumerate}

The next sections carry out these two reductions.


% ============================================================
% 5. Bracket-vector lower bound
% ============================================================

\section{Bracket-vector lower bound}
\label{sec:bracket-vector}

The purpose of this section is to prove that the reflected bracket vector does
not degenerate at the critical line.  This is the main structural reason the
inner-strip argument is different from the outer-strip argument.

The one-channel bracket
\[
        \Beff(\sigma,\tau)
\]
has a lower bound away from \(\sigma=1/2\), but that lower bound degenerates as
\(\sigma\to1/2\).  The two-channel bracket vector keeps both the value and the
transverse derivative at the critical line.  These two quantities cannot vanish
simultaneously.

Recall the effective bracket
\[
\begin{aligned}
        \Beff(\sigma,\tau)
        &=
        2^\sigma e^{\ii L}
        \left(
        \frac{12-\sigma}{48}
        +
        \frac{\ii X}{4}
        \right)
\\
        &\quad
        -
        2^{-\sigma}e^{-\ii L}
        \left(
        \frac{11+\sigma}{24}
        -
        \frac{\ii X}{2}
        \right),
\end{aligned}
\]
where
\[
        L=\tau\log2,
        \qquad
        X=\Xi_{\rm end}(\tau)\in\mathbb R.
\]
In this section \(L\) is treated as an arbitrary real phase and \(X\) as an
arbitrary real parameter.

\subsection{Critical-line bracket value and derivative}

Define the critical-line bracket value
\[
        B_0(L,X)
        :=
        \Beff\left(\frac12,\tau\right),
\]
and the critical-line transverse derivative
\[
        B_1(L,X)
        :=
        \partial_\sigma\Beff(\sigma,\tau)\bigg|_{\sigma=1/2}.
\]

\begin{lemma}[Critical-line bracket value]
\label{lem:critical-bracket-value}
For every real \(L\) and every real \(X\),
\[
        B_0(L,X)
        =
        \frac{\ii\sqrt2}{48}
        \left(
        24X\cos L+23\sin L
        \right).
\]
In particular, \(B_0(L,X)\) is purely imaginary.
\end{lemma}

\begin{proof}
Substitute \(\sigma=1/2\) into the bracket formula.  Since
\[
        2^{1/2}=\sqrt2,
        \qquad
        2^{-1/2}=\frac1{\sqrt2},
\]
we get
\[
\begin{aligned}
        B_0(L,X)
        &=
        \sqrt2 e^{\ii L}
        \left(
        \frac{23}{96}
        +
        \frac{\ii X}{4}
        \right)
        -
        \frac1{\sqrt2}e^{-\ii L}
        \left(
        \frac{23}{48}
        -
        \frac{\ii X}{2}
        \right).
\end{aligned}
\]
Putting both terms over the same denominator gives
\[
        B_0(L,X)
        =
        \frac{\sqrt2}{96}
        \left[
        e^{\ii L}(23+24\ii X)
        -
        e^{-\ii L}(46-48\ii X)
        \right].
\]
Expanding into sine and cosine terms and simplifying gives
\[
        B_0(L,X)
        =
        \frac{\ii\sqrt2}{48}
        \left(
        24X\cos L+23\sin L
        \right).
\]
\end{proof}

\begin{lemma}[Critical-line bracket derivative]
\label{lem:critical-bracket-derivative}
For every real \(L\) and every real \(X\),
\[
        B_1(L,X)
        =
        \frac{\sqrt2}{48}
        \left[
        (23\log2-2)\cos L
        -
        24X\log2\sin L
        \right].
\]
In particular, \(B_1(L,X)\) is real.
\end{lemma}

\begin{proof}
Differentiate the bracket formula with respect to \(\sigma\).  The first
bracket component gives
\[
\begin{aligned}
        \partial_\sigma
        \left[
        2^\sigma e^{\ii L}
        \left(
        \frac{12-\sigma}{48}
        +
        \frac{\ii X}{4}
        \right)
        \right]
        &=
        2^\sigma e^{\ii L}
        \left[
        \log2
        \left(
        \frac{12-\sigma}{48}
        +
        \frac{\ii X}{4}
        \right)
        -
        \frac1{48}
        \right].
\end{aligned}
\]
The second bracket component gives
\[
\begin{aligned}
        \partial_\sigma
        \left[
        -
        2^{-\sigma}e^{-\ii L}
        \left(
        \frac{11+\sigma}{24}
        -
        \frac{\ii X}{2}
        \right)
        \right]
        &=
        2^{-\sigma}e^{-\ii L}
        \left[
        \log2
        \left(
        \frac{11+\sigma}{24}
        -
        \frac{\ii X}{2}
        \right)
        -
        \frac1{24}
        \right].
\end{aligned}
\]
Substituting \(\sigma=1/2\), collecting the \(e^{\ii L}\) and \(e^{-\ii L}\)
terms, and simplifying yields
\[
        B_1(L,X)
        =
        \frac{\sqrt2}{48}
        \left[
        (23\log2-2)\cos L
        -
        24X\log2\sin L
        \right].
\]
The displayed expression is real because \(L\), \(X\), and \(\log2\) are real.
\end{proof}

\subsection{Critical-line vector nondegeneracy}

We now prove that \(B_0\) and \(B_1\) cannot vanish simultaneously.  More
precisely, their combined norm has a uniform positive lower bound.

Let
\[
        \ell:=\log2,
        \qquad
        a:=23\ell-2.
\]
Since
\[
        0.69<\ell<0.70,
\]
we have
\[
        a>0.
\]

\begin{proposition}[Critical-line bracket-vector lower bound]
\label{prop:critical-vector-lower-bound}
For every real \(L\) and every real \(X\),
\[
        |B_0(L,X)|^2+|B_1(L,X)|^2
        \ge
        \left(
        \frac{\sqrt2}{48}(23\log2-2)
        \right)^2.
\]
In particular,
\[
        |B_0(L,X)|^2+|B_1(L,X)|^2>0.
\]
\end{proposition}

\begin{proof}
By \cref{lem:critical-bracket-value,lem:critical-bracket-derivative},
\[
        |B_0(L,X)|^2+|B_1(L,X)|^2
        =
        \frac{2}{48^2}
        \left[
        A(L,X)^2+D(L,X)^2
        \right],
\]
where
\[
        A(L,X):=24X\cos L+23\sin L,
\]
and
\[
        D(L,X):=(23\ell-2)\cos L-24X\ell\sin L.
\]
Set
\[
        c:=\cos L,
        \qquad
        s:=\sin L,
        \qquad
        a:=23\ell-2.
\]
Then
\[
        A=24Xc+23s,
        \qquad
        D=ac-24X\ell s.
\]

For fixed \(L\), the quantity
\[
        A^2+D^2
\]
is a nonnegative quadratic polynomial in \(X\).  Its minimum over real \(X\)
is obtained by elementary projection.  A direct calculation gives
\[
        \min_{X\in\mathbb R}(A^2+D^2)
        =
        \frac{(23\ell-2\cos^2 L)^2}
        {\ell^2+(1-\ell^2)\cos^2 L}.
\]
Writing
\[
        t:=\cos^2 L,
        \qquad
        0\le t\le1,
\]
this minimum becomes
\[
        \frac{(23\ell-2t)^2}{\ell^2+(1-\ell^2)t}.
\]
We claim that
\[
        \frac{(23\ell-2t)^2}{\ell^2+(1-\ell^2)t}
        \ge
        (23\ell-2)^2
        =
        a^2
\]
for \(0\le t\le1\).

Indeed, since the denominator
\[
        \ell^2+(1-\ell^2)t
\]
is positive, it is enough to show
\[
        (23\ell-2t)^2
        -
        (23\ell-2)^2
        \left[
        \ell^2+(1-\ell^2)t
        \right]
        \ge0.
\]
The left-hand side factors as
\[
        (t-1)
        \left(
        529\ell^4
        -
        92\ell^3
        -
        525\ell^2
        +
        4t
        \right).
\]
Since \(0\le t\le1\), we have
\[
        t-1\le0.
\]
Also, since \(0\le t\le1\), \(0.69<\ell<0.70\), and all coefficients are
explicit, we have the rigorous upper bound
\[
\begin{aligned}
        529\ell^4
        -
        92\ell^3
        -
        525\ell^2
        +
        4t
        &\le
        529(0.70)^4
        -
        92(0.69)^3
        -
        525(0.69)^2
        +
        4
\\
        &<
        0.
\end{aligned}
\]
Hence
\[
        t-1\le0
\]
and the second factor is strictly negative.  Their product is therefore
nonnegative.  Therefore
\[
        A^2+D^2\ge a^2.
\]
Consequently,
\[
        |B_0(L,X)|^2+|B_1(L,X)|^2
        \ge
        \frac{2}{48^2}a^2
        =
        \left(
        \frac{\sqrt2}{48}(23\log2-2)
        \right)^2.
\]
This proves the proposition.
\end{proof}

\begin{remark}[Interpretation]
The outer-strip one-channel lower bound degenerates because
\[
        |\Beff(\sigma,\tau)|
\]
can become small as \(\sigma\to1/2\).  The proposition says that the pair
\[
        \left(
        \Beff\left(\frac12,\tau\right),
        \partial_\sigma\Beff\left(\frac12,\tau\right)
        \right)
\]
does not degenerate.  Thus a reflected zero pair would have to cancel not only
the bracket value at the critical line but also its transverse slope.  The
bracket algebra forbids this simultaneous cancellation.
\end{remark}

\subsection{Finite-width bracket certificate}

The previous proposition proves the limiting lower bound at \(u=0\).  The
inner-strip detector uses the finite-width bracket vector
\[
        \Bvec(u,\tau)
        =
        \left(
        B_{\rm ev}(u,\tau),
        B_{\rm od}(u,\tau)
        \right),
\]
where
\[
        0<u\le0.01.
\]
As \(u\to0\),
\[
        B_{\rm ev}(u,\tau)
        \to
        B_0(L,X),
\]
and
\[
        B_{\rm od}(u,\tau)
        \to
        -B_1(L,X).
\]
The critical-line lower bound is approximately
\[
        \frac{\sqrt2}{48}(23\log2-2)
        =
        0.4107814619\ldots.
\]
For the public theorem, we use the deliberately smaller finite-width target
\[
        \gev=0.2.
\]

The finite-width inequality
\[
        \|\Bvec(u,\tau)\|\ge0.2
\]
for
\[
        0<u\le0.01
\]
is certified by a finite real interval calculation as follows.  The phase
\(L\) is reduced modulo \(2\pi\).  For fixed \(u\) and \(L\), the squared norm
\[
        \|\Bvec(u,\tau)\|^2
\]
is an explicit quadratic polynomial in the real parameter
\[
        X=\Xi_{\rm end}(\tau).
\]
Thus the dependence on \(X\in\mathbb R\) is removed by minimizing this
quadratic analytically over \(X\).  The remaining verification is a compact
two-variable interval problem in
\[
        0\le u\le0.01,
        \qquad
        0\le L\le2\pi.
\]
The public certificate proves the conservative bound
\[
        \|\Bvec(u,\tau)\|\ge0.2.
\]

\begin{assumption}[Finite inner bracket-vector certificate]
\label{ass:finite-bracket-vector-certificate}
The public file
\[
        \texttt{inner\_bracket\_vector\_certificate.json}
\]
has status
\[
        \texttt{proved}
\]
and certifies
\[
        \|\Bvec(u,\tau)\|\ge\gev
\]
for
\[
        0<u\le0.01,
        \qquad
        \tau\ge T_0,
\]
with
\[
        \gev=0.2.
\]
\end{assumption}

\begin{remark}[Why a finite certificate is acceptable here]
The limiting lower bound at \(u=0\) is analytic.  Passing from \(u=0\) to
\(0<u\le0.01\) is a finite compact verification problem for an explicit
elementary expression.  No zeta-function estimates, no stationary phase, and
no residual analysis enter this bracket-vector certificate.  It is a small
real interval problem, separate from the much larger residual certificate.
\end{remark}

\begin{proposition}[Certified inner bracket-vector lower bound]
\label{prop:certified-inner-bracket-vector}
Assume \cref{ass:finite-bracket-vector-certificate}.  Then, for every
\[
        0<u\le0.01,
        \qquad
        \tau\ge T_0,
\]
one has
\[
        \|\Bvec(u,\tau)\|\ge0.2.
\]
\end{proposition}

\begin{proof}
This is exactly the statement certified by
\cref{ass:finite-bracket-vector-certificate}.
\end{proof}


% ============================================================
% 6. Residual-vector bound from value and derivative certificates
% ============================================================

\section{Residual-vector bound}
\label{sec:residual-vector}

We now turn from the bracket vector to the residual vector.  This is the place
where the inner-strip argument genuinely differs from the outer-strip
argument.

In the outer strip, it is enough to certify the residual itself:
\[
        |\Rtot(s)|
        \le
        C
        \frac{\log\tau}{\tau^{3/2}}.
\]
For the inner reflected detector, this controls the even channel but not the
divided odd channel.  The odd channel contains the quotient
\[
        \frac{\Rtot(s_-)-\Rtot(s_+)}{2u}.
\]
A uniform bound for this expression as \(u\to0\) requires a bound for
\[
        \partial_\sigma\Rtot(\sigma+\ii\tau).
\]
Thus the inner-strip certificate must include both ordinary residual constants
and \(\sigma\)-derivative residual constants.

\subsection{Sector-level assumptions}

For clarity we restate the sector certificate interface in the form used in
this section.  The total residual is decomposed as
\[
        \Rtot
        =
        \sum_{\nu\in\mathcal S}
        \mathcal R_\nu,
\]
where
\[
        \mathcal S
        :=
        \{
        {\rm band},
        {\rm bd},
        {\rm far},
        {\rm end},
        {\rm nonstat},
        {\rm floor},
        {\rm core}
        \}.
\]

\begin{assumption}[Inner residual value and derivative certificates]
\label{ass:inner-residual-certificates}
For each
\[
        \nu\in\mathcal S,
\]
there exist constants
\[
        C_{\nu,0}\ge0,
        \qquad
        C_{\nu,1}\ge0,
\]
such that, for every
\[
        \left|\sigma-\frac12\right|\le0.01,
        \qquad
        \tau\ge T_0,
\]
one has
\[
        |\mathcal R_\nu(\sigma+\ii\tau)|
        \le
        C_{\nu,0}
        \frac{\log\tau}{\tau^{3/2}},
\]
and
\[
        |\partial_\sigma\mathcal R_\nu(\sigma+\ii\tau)|
        \le
        C_{\nu,1}
        \frac{\log\tau}{\tau^{3/2}}.
\]
\end{assumption}

Define
\[
        \Czero
        :=
        \sum_{\nu\in\mathcal S}
        C_{\nu,0},
\]
and
\[
        \Cone
        :=
        \sum_{\nu\in\mathcal S}
        C_{\nu,1}.
\]
Finally define
\[
        \Cinner
        :=
        \sqrt{\Czero^2+\Cone^2}.
\]

\begin{remark}[Why the constants are summed before taking the vector norm]
The residual vector has two components.  The even component is bounded by the
sum of the ordinary sector constants, and the divided odd component is bounded
by the sum of the derivative sector constants.  Only after the two scalar
bounds are obtained do we take the Euclidean vector norm.  This is why
\[
        \Cinner=\sqrt{\Czero^2+\Cone^2}
\]
is the correct inner residual constant.
\end{remark}

\subsection{Even residual channel}

Recall
\[
        R_{\rm ev}(u,\tau)
        =
        \frac{\Rtot(s_-)+\Rtot(s_+)}2,
\]
where
\[
        s_-=\frac12-u+\ii\tau,
        \qquad
        s_+=\frac12+u+\ii\tau.
\]

\begin{lemma}[Even residual bound]
\label{lem:even-residual-bound}
Assume \cref{ass:inner-residual-certificates}.  For
\[
        0<u\le0.01,
        \qquad
        \tau\ge T_0,
\]
one has
\[
        |R_{\rm ev}(u,\tau)|
        \le
        \Czero
        \frac{\log\tau}{\tau^{3/2}}.
\]
\end{lemma}

\begin{proof}
By the triangle inequality,
\[
        |R_{\rm ev}(u,\tau)|
        \le
        \frac12
        |\Rtot(s_-)|
        +
        \frac12
        |\Rtot(s_+)|.
\]
Since
\[
        \Rtot=\sum_{\nu\in\mathcal S}\mathcal R_\nu,
\]
we have
\[
        |\Rtot(s_\pm)|
        \le
        \sum_{\nu\in\mathcal S}
        |\mathcal R_\nu(s_\pm)|.
\]
Both \(s_-\) and \(s_+\) lie in the certified inner strip
\[
        \left|\sigma-\frac12\right|\le0.01.
\]
Therefore, by \cref{ass:inner-residual-certificates},
\[
        |\mathcal R_\nu(s_\pm)|
        \le
        C_{\nu,0}
        \frac{\log\tau}{\tau^{3/2}}.
\]
Summing over sectors gives
\[
        |\Rtot(s_\pm)|
        \le
        \Czero
        \frac{\log\tau}{\tau^{3/2}}.
\]
Hence
\[
        |R_{\rm ev}(u,\tau)|
        \le
        \Czero
        \frac{\log\tau}{\tau^{3/2}}.
\]
\end{proof}

\subsection{Divided odd residual channel}

The divided odd residual channel is
\[
        R_{\rm od}(u,\tau)
        =
        \frac{\Rtot(s_-)-\Rtot(s_+)}{2u}.
\]
The point of certifying \(\partial_\sigma\Rtot\) is exactly to control this
quantity uniformly as \(u\to0\).

\begin{lemma}[Divided odd residual bound]
\label{lem:odd-residual-bound}
Assume \cref{ass:inner-residual-certificates}.  For
\[
        0<u\le0.01,
        \qquad
        \tau\ge T_0,
\]
one has
\[
        |R_{\rm od}(u,\tau)|
        \le
        \Cone
        \frac{\log\tau}{\tau^{3/2}}.
\]
\end{lemma}

\begin{proof}
For fixed \(\tau\), define
\[
        F(\sigma):=\Rtot(\sigma+\ii\tau).
\]
Then
\[
        R_{\rm od}(u,\tau)
        =
        \frac{
        F(\frac12-u)-F(\frac12+u)
        }{2u}.
\]
By the mean value theorem applied to the real variable \(\sigma\),
\[
        |R_{\rm od}(u,\tau)|
        \le
        \sup_{\left|\sigma-\frac12\right|\le u}
        |\partial_\sigma\Rtot(\sigma+\ii\tau)|.
\]
Since \(u\le0.01\), the supremum is taken inside the certified inner strip
\[
        \left|\sigma-\frac12\right|\le0.01.
\]
Using
\[
        \Rtot=\sum_{\nu\in\mathcal S}\mathcal R_\nu,
\]
we have
\[
        |\partial_\sigma\Rtot(\sigma+\ii\tau)|
        \le
        \sum_{\nu\in\mathcal S}
        |\partial_\sigma\mathcal R_\nu(\sigma+\ii\tau)|.
\]
By \cref{ass:inner-residual-certificates},
\[
        |\partial_\sigma\mathcal R_\nu(\sigma+\ii\tau)|
        \le
        C_{\nu,1}
        \frac{\log\tau}{\tau^{3/2}}.
\]
Therefore
\[
        |\partial_\sigma\Rtot(\sigma+\ii\tau)|
        \le
        \Cone
        \frac{\log\tau}{\tau^{3/2}}.
\]
This proves
\[
        |R_{\rm od}(u,\tau)|
        \le
        \Cone
        \frac{\log\tau}{\tau^{3/2}}.
\]
\end{proof}

\begin{remark}[No loss as \(u\to0\)]
The divided odd residual channel contains \(1/u\), but the mean value theorem
removes this apparent singularity.  The price is exactly the
\(\sigma\)-derivative residual certificate.  This is why the inner-strip
certificate is finite but strictly stronger than the outer-strip certificate.
\end{remark}

\subsection{Residual-vector bound}

We now combine the two scalar residual estimates.

\begin{proposition}[Two-channel residual-vector bound]
\label{prop:residual-vector-bound}
Assume \cref{ass:inner-residual-certificates}.  For
\[
        0<u\le0.01,
        \qquad
        \tau\ge T_0,
\]
the residual vector
\[
        \Rvec(u,\tau)
        =
        \left(
        R_{\rm ev}(u,\tau),
        R_{\rm od}(u,\tau)
        \right)
\]
satisfies
\[
        \|\Rvec(u,\tau)\|
        \le
        \Cinner
        \frac{\log\tau}{\tau^{3/2}},
\]
where
\[
        \Cinner
        =
        \sqrt{\Czero^2+\Cone^2}.
\]
\end{proposition}

\begin{proof}
By \cref{lem:even-residual-bound},
\[
        |R_{\rm ev}(u,\tau)|
        \le
        \Czero
        \frac{\log\tau}{\tau^{3/2}}.
\]
By \cref{lem:odd-residual-bound},
\[
        |R_{\rm od}(u,\tau)|
        \le
        \Cone
        \frac{\log\tau}{\tau^{3/2}}.
\]
Therefore, using the Euclidean norm on \(\mathbb C^2\),
\[
\begin{aligned}
        \|\Rvec(u,\tau)\|
        &=
        \left(
        |R_{\rm ev}(u,\tau)|^2
        +
        |R_{\rm od}(u,\tau)|^2
        \right)^{1/2}
\\
        &\le
        \left(
        \Czero^2+\Cone^2
        \right)^{1/2}
        \frac{\log\tau}{\tau^{3/2}}
\\
        &=
        \Cinner
        \frac{\log\tau}{\tau^{3/2}}.
\end{aligned}
\]
\end{proof}

\subsection{Certificate form of the residual-vector theorem}

For publication and reproducibility, the constants in
\cref{prop:residual-vector-bound} are not to be copied by hand from separate
calculations.  They must be produced by a strict merge of the sector
certificates.

\begin{assumption}[Final inner residual certificate]
\label{ass:final-inner-residual-certificate}
The file
\[
        \texttt{INNER\_STRIP\_CERTIFICATE\_SAFE.json}
\]
contains, for every
\[
        \nu\in\mathcal S,
\]
a dependency file with status
\[
        \texttt{proved}
\]
and with fields
\[
        \texttt{C\_value\_upper\_safe},
        \qquad
        \texttt{C\_sigma\_derivative\_upper\_safe}.
\]
The final file has status
\[
        \texttt{inner\_full\_certificate}
\]
and records the merged values
\[
        \Czero,
        \qquad
        \Cone,
        \qquad
        \Cinner.
\]
\end{assumption}

\begin{corollary}[Certified residual-vector estimate]
\label{cor:certified-residual-vector-estimate}
Assume \cref{ass:final-inner-residual-certificate}.  Then the residual vector
satisfies
\[
        \|\Rvec(u,\tau)\|
        \le
        \Cinner
        \frac{\log\tau}{\tau^{3/2}}
\]
for every
\[
        0<u\le0.01,
        \qquad
        \tau\ge T_0.
\]
\end{corollary}

\begin{proof}
The final certificate supplies the sector constants satisfying
\cref{ass:inner-residual-certificates}.  The result then follows from
\cref{prop:residual-vector-bound}.
\end{proof}

\begin{remark}[What the final certificate must refuse]
A final certificate is not acceptable if any sector lacks either its ordinary
value constant or its \(\sigma\)-derivative constant.  In particular, a file
which certifies only
\[
        |\mathcal R_\nu(s)|
\]
but not
\[
        |\partial_\sigma\mathcal R_\nu(s)|
\]
does not certify the inner-strip theorem.  The divided odd channel requires the
derivative data.
\end{remark}

% ============================================================
% 7. Finite certifiability of the derivative residual sectors
% ============================================================

\section{Finite certifiability of the derivative residual sectors}
\label{sec:sector-certifiability}

The previous section showed that the inner-strip residual vector is controlled
once every residual sector is certified both in value and in
\(\sigma\)-derivative.  We now explain why these derivative certificates are
finite in exactly the same sense as the outer-strip residual certificate.

The point of this section is not to reproduce every numerical enclosure.
Those belong to the public certificate files.  The point is to show that the
\(\sigma\)-derivative operation does not introduce any new analytic
non-compactness, any new stationary phase family, or any untracked sector.  It
only differentiates explicit amplitudes, explicit algebraic coefficients, and
the explicit quotient projection.

Thus the inner-strip derivative certificate is a finite enhancement of the
outer-strip certificate, not a new asymptotic theory.

\subsection{General derivative principle}

Each residual sector in the quotient expansion has the schematic form
\[
        \mathcal R_\nu(s)
        =
        \Pi_Q(s)\,
        \mathcal A_\nu(s)
\]
or a finite sum of such terms, where
\[
        \Pi_Q(s)
        =
        \frac{\ii\tau}{(1-s)(s+1)}
\]
is the quotient projection and \(\mathcal A_\nu(s)\) is a sector amplitude,
integral, or grouped mode sum.

The \(\sigma\)-derivative is therefore
\[
        \partial_\sigma \mathcal R_\nu(s)
        =
        (\partial_\sigma\Pi_Q(s))\mathcal A_\nu(s)
        +
        \Pi_Q(s)\partial_\sigma\mathcal A_\nu(s),
\]
with
\[
        \partial_\sigma\Pi_Q(s)
        =
        \Pi_Q(s)
        \left(
        \frac1{1-s}
        -
        \frac1{s+1}
        \right).
\]
Equivalently,
\[
        \partial_\sigma\Pi_Q(s)
        =
        \frac{2s\,\ii\tau}{[(1-s)(s+1)]^2}.
\]

\begin{remark}[Projection derivative is mandatory]
A derivative certificate which differentiates the sector amplitudes but omits
\[
        \partial_\sigma\Pi_Q(s)
\]
does not certify the inner-strip theorem.  The quotient projection depends on
\(\sigma\), and its derivative is part of the residual derivative.
\end{remark}

On the inner strip
\[
        \left|\sigma-\frac12\right|\le0.01,
        \qquad
        \tau\ge10,
\]
the projection and its \(\sigma\)-derivative are elementary rational functions
of \(s\) and are uniformly bounded by explicit powers of \(1/\tau\).  Thus the
projection derivative introduces no new analytic difficulty.

\subsection{Algebraic coefficient derivatives}

The two main \(J\)-side coefficients are
\[
        s(s-2)
\]
and
\[
        (1-s)(s+1).
\]
Their \(\sigma\)-derivatives are
\[
        \partial_\sigma[s(s-2)]
        =
        2s-2,
\]
and
\[
        \partial_\sigma[(1-s)(s+1)]
        =
        -2s.
\]
These are explicit algebraic factors.  On the inner strip they have size
\(O(\tau)\), whereas the original quadratic coefficients have size
\(O(\tau^2)\).  Therefore differentiating the coefficients never worsens the
asymptotic order of any sector after the same quotient normalization is
applied.

\begin{remark}[No hidden coefficient sector]
The coefficient derivatives are assigned to the same sector as the
undifferentiated coefficient.  They are not a new residual type and do not
create an eighth sector.
\end{remark}

\subsection{Compact-boundary sector}

The compact-boundary sector is the delicate sector in the outer-strip
certificate.  It is also the most important sector for the derivative
certificate.

In the compact-boundary scaling, the relevant amplitudes are built from factors
of the form
\[
        (2+wy)^{-\sigma-2}
\]
and
\[
        (2+wy)^{\sigma-3},
\]
together with their frozen endpoint counterparts
\[
        2^{-\sigma-2},
        \qquad
        2^{\sigma-3}.
\]
Their \(\sigma\)-derivatives are explicit:
\[
        \partial_\sigma(2+wy)^{-\sigma-2}
        =
        -\log(2+wy)(2+wy)^{-\sigma-2},
\]
\[
        \partial_\sigma(2+wy)^{\sigma-3}
        =
        \log(2+wy)(2+wy)^{\sigma-3},
\]
and
\[
        \partial_\sigma 2^{-\sigma-2}
        =
        -\log2\,2^{-\sigma-2},
\]
\[
        \partial_\sigma 2^{\sigma-3}
        =
        \log2\,2^{\sigma-3}.
\]

In the compact-boundary certificate the cutoff satisfies
\[
        0\le y\le2Y,
        \qquad
        Y=128.
\]
Therefore
\[
        0\le y\le256.
\]
For
\[
        0\le w\le\frac1{\sqrt{10}},
\]
the logarithmic factors
\[
        \log(2+wy)
\]
are bounded on the compact certificate domain.  Hence differentiating in
\(\sigma\) only inserts bounded logarithmic amplitude factors inside the same
finite collection of Arb boxes.

The compact-boundary derivative verifier must therefore enclose, on every box,
both the original compact-boundary integrand and its \(\sigma\)-derivative.
For a grouped contribution of the form
\[
        G_\nu(s)
        =
        c_+(s)\Delta_+(s)
        +
        c_-(s)\Delta_-(s),
\]
the derivative is
\[
        \partial_\sigma G_\nu(s)
        =
        c_+'(s)\Delta_+(s)
        +
        c_+(s)\partial_\sigma\Delta_+(s)
        +
        c_-'(s)\Delta_-(s)
        +
        c_-(s)\partial_\sigma\Delta_-(s).
\]
If this grouped contribution is then multiplied by the quotient projection,
one must also apply
\[
        \partial_\sigma(\Pi_Q G_\nu)
        =
        (\partial_\sigma\Pi_Q)G_\nu
        +
        \Pi_Q\partial_\sigma G_\nu.
\]

\begin{proposition}[Finite certifiability of the compact-boundary derivative]
\label{prop:bd-derivative-certifiable}
The compact-boundary sector admits a finite \(\sigma\)-derivative certificate
in the inner-strip norm
\[
        |\partial_\sigma\mathcal R_{\rm bd}(s)|
        \le
        C_{{\rm bd},1}
        \frac{\log\tau}{\tau^{3/2}}.
\]
\end{proposition}

\begin{proof}
The compact-boundary sector is already reduced to finitely many compact boxes
in the variables
\[
        \sigma,\quad w,\quad r,\quad y,\quad \lambda.
\]
Differentiating with respect to \(\sigma\) leaves the phases unchanged and
only differentiates explicit amplitudes, explicit algebraic coefficients, and
the quotient projection.  The amplitude derivatives insert the bounded factors
\[
        \log(2+wy)
        \quad\text{and}\quad
        \log2
\]
on the compact support \(0\le y\le256\).  The coefficient and projection
derivatives are explicit rational functions of \(s\).  Therefore the same
finite box subdivision certifies the derivative integrands by Arb ball
arithmetic.  Summing the resulting box enclosures gives a finite constant
\(C_{{\rm bd},1}\).
\end{proof}

\subsection{Compact-complement sectors}

The compact complement is split into the same three hooks as in the
outer-strip certificate:
\[
        C_{\chi^c}
        =
        C_{\chi^c,{\rm tail}}
        +
        C_{\chi^c,{\rm nonstat}}
        +
        C_{\chi^c,{\rm stationary}}.
\]
For the inner strip, each hook receives a derivative companion.

\subsubsection*{Tail hook}

The compact-complement tail hook is controlled by derivative envelopes followed
by integrations by parts.  The derivative certificate requires envelopes not
only for the original amplitude derivatives but also for their
\(\sigma\)-derivatives.  Schematically, if the original tail hook certifies
\[
        |H'(y)|\le M_1,
        \qquad
        |H''(y)|\le M_2,
\]
then the derivative hook certifies
\[
        |\partial_\sigma H'(y)|\le M_{1,\sigma},
        \qquad
        |\partial_\sigma H''(y)|\le M_{2,\sigma}.
\]
The same integrations by parts then produce the derivative tail constant
\[
        C_{\chi^c,{\rm tail},1}.
\]

\subsubsection*{Nonstationary hook}

The nonstationary hook is based on a lower bound for the phase derivative.  The
phase derivative is independent of \(\sigma\); only amplitudes and coefficients
are differentiated.  Hence the same nonstationary box partition and the same
phase-separation condition apply to the derivative integrands.  The
integration-by-parts bound is repeated with the \(\sigma\)-differentiated
amplitudes.

\subsubsection*{Stationary-overlap hook}

The stationary-overlap hook is compact and is certified by direct Arb
enclosure.  Its derivative version is obtained by enclosing the
\(\sigma\)-differentiated integrand on the same compact boxes.

\begin{proposition}[Finite certifiability of compact-complement derivatives]
\label{prop:chic-derivative-certifiable}
Each compact-complement hook admits a finite \(\sigma\)-derivative certificate:
\[
        C_{\chi^c,{\rm tail},1},
        \qquad
        C_{\chi^c,{\rm nonstat},1},
        \qquad
        C_{\chi^c,{\rm stationary},1}.
\]
Consequently the full compact-complement derivative contribution is finitely
certifiable.
\end{proposition}

\begin{proof}
The tail hook uses derivative-envelope bounds; differentiating with respect to
\(\sigma\) replaces the original envelopes by envelopes for the
\(\sigma\)-differentiated amplitudes.  The nonstationary hook uses the same
phase-separation boxes because the phases do not depend on \(\sigma\).  The
stationary-overlap hook is compact and is directly enclosed after
differentiation.  Each operation is performed on a finite box subdivision, and
all differentiated amplitudes are explicit elementary functions.  Therefore
all three derivative constants are finite and certifiable.
\end{proof}

\subsection{Band sector}

The band sector comes from the quotient stationary range
\[
        \frac{\tau}{4\pi}<m\le\frac{\tau}{2\pi}.
\]
The relevant sums contain powers such as
\[
        m^{-s},
        \qquad
        m^{-(1-s)}.
\]
Differentiating in \(\sigma\) gives
\[
        \partial_\sigma m^{-s}
        =
        -(\log m)m^{-s},
\]
and
\[
        \partial_\sigma m^{-(1-s)}
        =
        (\log m)m^{-(1-s)}.
\]
On the band,
\[
        m\asymp \tau,
\]
so
\[
        \log m\le C\log\tau
\]
for an explicit constant \(C\).  The Poisson endpoint expansion and the
nonstationary dual-frequency estimates therefore remain finite after
differentiation, with derivative amplitudes containing at most logarithmic
factors.

\begin{proposition}[Finite certifiability of the band derivative]
\label{prop:band-derivative-certifiable}
The band residual admits a finite \(\sigma\)-derivative certificate
\[
        |\partial_\sigma\mathcal R_{\rm band}(s)|
        \le
        C_{{\rm band},1}
        \frac{\log\tau}{\tau^{3/2}}
\]
on the certified inner strip.
\end{proposition}

\begin{proof}
The differentiated band sums are obtained from the original band sums by
inserting factors \(\pm\log m\).  In the band range \(m\asymp\tau\), these
factors are bounded by explicit multiples of \(\log\tau\).  The same Poisson
endpoint decomposition applies.  The endpoint stationary contributions are
differentiated explicitly, and the nonstationary dual-frequency tails are
bounded by the same integration-by-parts scheme with differentiated
amplitudes.  Thus the band derivative is controlled in the same residual norm,
with a finite constant \(C_{{\rm band},1}\).
\end{proof}

\subsection{Endpoint sector}

The endpoint sector contains the regularized endpoint scalar
\[
        X=\Xi_{\rm end}(\tau).
\]
This scalar depends on \(\tau\), but not on \(\sigma\).  Therefore
\(\partial_\sigma\) acts only on the explicit powers and coefficients in the
endpoint bracket and endpoint residual.

The relevant differentiated factors are combinations of
\[
        2^\sigma,
        \qquad
        2^{-\sigma},
        \qquad
        12-\sigma,
        \qquad
        11+\sigma,
\]
whose derivatives are elementary:
\[
        \partial_\sigma 2^\sigma=(\log2)2^\sigma,
\]
\[
        \partial_\sigma 2^{-\sigma}=-(\log2)2^{-\sigma},
\]
\[
        \partial_\sigma(12-\sigma)=-1,
\]
and
\[
        \partial_\sigma(11+\sigma)=1.
\]

\begin{proposition}[Finite certifiability of the endpoint derivative]
\label{prop:endpoint-derivative-certifiable}
The endpoint residual admits a finite \(\sigma\)-derivative certificate
\[
        |\partial_\sigma\mathcal R_{\rm end}(s)|
        \le
        C_{{\rm end},1}
        \frac{\log\tau}{\tau^{3/2}}.
\]
\end{proposition}

\begin{proof}
The endpoint residual is an explicit finite combination of the real scalar
\(\Xi_{\rm end}(\tau)\), powers \(2^{\pm\sigma}\), and algebraic coefficients
in \(\sigma\).  Since \(\Xi_{\rm end}\) is independent of \(\sigma\), the
\(\sigma\)-derivative is obtained by differentiating only elementary factors.
The resulting expression is explicit and is bounded by the same interval
method used for the endpoint sector itself.
\end{proof}

\subsection{Far stationary sector}

The far stationary sector is controlled away from the endpoint-stationary
transition.  Its amplitudes depend on \(\sigma\) through powers of the
stationary point
\[
        x_{m,*}
        =
        \frac{\tau}{2\pi m}
\]
and through explicit algebraic coefficients.  Differentiating those powers
inserts logarithmic factors such as
\[
        \log x_{m,*}.
\]
In the far stationary range, \(x_{m,*}\) stays in the same controlled domain as
in the original stationary analysis.  The differentiated stationary remainder
therefore has the same summability structure, with explicit logarithmic
weights.

\begin{proposition}[Finite certifiability of the far stationary derivative]
\label{prop:far-derivative-certifiable}
The far stationary residual admits a finite \(\sigma\)-derivative certificate
\[
        |\partial_\sigma\mathcal R_{\rm far}(s)|
        \le
        C_{{\rm far},1}
        \frac{\log\tau}{\tau^{3/2}}.
\]
\end{proposition}

\begin{proof}
The far stationary expansion is uniform away from the boundary transition.
Differentiating with respect to \(\sigma\) differentiates only amplitudes and
coefficients; the stationary phase location and curvature are independent of
\(\sigma\).  The differentiated amplitudes contain logarithmic factors
\(\log x_{m,*}\), which are explicitly bounded on the far stationary boxes.
The same summation and interval enclosure used for the undifferentiated far
sector therefore gives a finite derivative constant.
\end{proof}

\subsection{Always-nonstationary sector}

The always-nonstationary sector is controlled by integration by parts in
regions where the phase derivative is separated from zero.  The phases do not
depend on \(\sigma\).  Therefore the nonstationary separation is unchanged
under \(\partial_\sigma\).

Differentiating the sector only differentiates amplitudes, coefficients, and
possibly the quotient projection.  The differentiated amplitudes contain
bounded logarithmic factors on each certified box.

\begin{proposition}[Finite certifiability of the nonstationary derivative]
\label{prop:nonstat-derivative-certifiable}
The always-nonstationary residual admits a finite \(\sigma\)-derivative
certificate
\[
        |\partial_\sigma\mathcal R_{\rm nonstat}(s)|
        \le
        C_{{\rm nonstat},1}
        \frac{\log\tau}{\tau^{3/2}}.
\]
\end{proposition}

\begin{proof}
The integration-by-parts denominators are phase derivatives and are independent
of \(\sigma\).  Hence the same lower bounds for the phase derivatives apply to
the differentiated sector.  The \(\sigma\)-derivative acts only on elementary
amplitudes, coefficients, and the projection.  These differentiated terms are
explicit and are enclosed on the same finite partition.  Thus the
nonstationary derivative sector is finitely certifiable.
\end{proof}

\subsection{Floor and core sectors}

The floor sector is eliminated by the same half-open convention as in the
outer-strip proof.  If
\[
        \mathcal R_{\rm floor}=0,
\]
then also
\[
        \partial_\sigma\mathcal R_{\rm floor}=0.
\]

The core sector consists of algebraic quotient-normalization corrections and
coefficient mismatches.  These are explicit functions of \(s\), \(\tau\), and
the quotient projection.  Differentiating the core sector gives explicit
rational functions and elementary factors, all bounded on the inner strip.

\begin{proposition}[Finite certifiability of floor and core derivatives]
\label{prop:floor-core-derivative-certifiable}
The floor and core residuals admit finite derivative certificates.  In
particular, if the floor sector is identically zero, then
\[
        C_{{\rm floor},1}=0.
\]
The core derivative satisfies
\[
        |\partial_\sigma\mathcal R_{\rm core}(s)|
        \le
        C_{{\rm core},1}
        \frac{\log\tau}{\tau^{3/2}}
\]
for a finite explicit constant \(C_{{\rm core},1}\).
\end{proposition}

\begin{proof}
The floor statement is immediate if the floor residual is identically zero.
The core sector is algebraic in the quotient coefficients, coefficient ratios,
and projection factors.  Its \(\sigma\)-derivative is obtained by differentiating
explicit rational functions of \(s\).  On
\[
        \left|\sigma-\frac12\right|\le0.01,
        \qquad
        \tau\ge10,
\]
these rational functions have explicit interval bounds.  Therefore the core
derivative is finitely certifiable.
\end{proof}

\subsection{Summary of derivative certifiability}

Combining the preceding propositions, every residual sector admits a
\(\sigma\)-derivative certificate in the inner-strip norm.  Therefore the only
remaining task is computational: produce the seven sector JSON files with both
ordinary and derivative constants, and merge them into the final public
certificate.

\begin{theorem}[Finite certifiability of the inner residual vector]
\label{thm:finite-certifiability-inner-residual}
The inner-strip residual-vector bound
\[
        \|\Rvec(u,\tau)\|
        \le
        \Cinner
        \frac{\log\tau}{\tau^{3/2}}
\]
is finitely certifiable from the seven sector certificates
\[
        \{
        {\rm band},
        {\rm bd},
        {\rm far},
        {\rm end},
        {\rm nonstat},
        {\rm floor},
        {\rm core}
        \},
\]
provided each sector certificate contains both a value constant and a
\(\sigma\)-derivative constant.
\end{theorem}

\begin{proof}
By
\cref{prop:bd-derivative-certifiable,prop:chic-derivative-certifiable,prop:band-derivative-certifiable,prop:endpoint-derivative-certifiable,prop:far-derivative-certifiable,prop:nonstat-derivative-certifiable,prop:floor-core-derivative-certifiable},
each sector derivative is finitely certifiable.  The ordinary sector residuals
are the same sector residuals as in the outer-strip certificate.  Once the
ordinary and derivative constants are available, the residual-vector bound is
exactly \cref{prop:residual-vector-bound}.
\end{proof}

\begin{remark}[What has been reduced to computation]
The analytic work in this section shows that no new infinite analytic problem
appears after differentiating in \(\sigma\).  The derivative certificate is a
finite interval-arithmetic problem on the same structural decomposition as the
outer-strip certificate.  The public proof is complete only after the final
certificate file records the corresponding constants.
\end{remark}



% ============================================================
% 8. Certified inner-strip exclusion theorem
% ============================================================

\section{Certified inner-strip exclusion}
\label{sec:inner-exclusion}

We now combine the reflected zero-pair identity, the two-channel
bracket-residual decomposition, the bracket-vector lower bound, and the
residual-vector estimate.

The proof is formally parallel to the outer-strip exclusion proof, but with two
important differences:

\begin{enumerate}[label=(\roman*)]
    \item the scalar bracket lower bound is replaced by the vector lower bound
    \[
            \|\Bvec(u,\tau)\|\ge \gev;
    \]

    \item the scalar residual bound is replaced by the residual-vector bound
    \[
            \|\Rvec(u,\tau)\|
            \le
            \Cinner
            \frac{\log\tau}{\tau^{3/2}}.
    \]
\end{enumerate}

The advantage of the two-channel formulation is that it remains meaningful as
\[
        u\to0.
\]

\begin{theorem}[Conditional certified inner-strip exclusion]
\label{thm:conditional-certified-inner-strip}
Assume the structural quotient expansion
\[
        \NQ(s)
        =
        \frac{\ii}{\tau}\Beff(\sigma,\tau)
        +
        \Rtot(s)
\]
from \cref{ass:inner-quotient-expansion}.  Assume also the finite bracket-vector
certificate of \cref{ass:finite-bracket-vector-certificate} and the final
inner residual certificate of \cref{ass:final-inner-residual-certificate}.

Let
\[
        0<u\le0.01,
        \qquad
        s_\pm=\frac12\pm u+\ii\tau.
\]
If \(T\ge T_0\) satisfies
\[
        \sqrt T
        >
        \frac{\Cinner}{\gev}
        \log T,
\]
then there is no reflected off-critical zero pair
\[
        s_-=\frac12-u+\ii\tau,
        \qquad
        s_+=\frac12+u+\ii\tau
\]
with
\[
        \tau\ge T.
\]
Equivalently, no nontrivial zero exists in the inner strip
\[
        0<\left|\Re s-\frac12\right|\le0.01,
        \qquad
        \Im s\ge T.
\]
\end{theorem}

\begin{proof}
Suppose, for contradiction, that
\[
        s_-=\frac12-u+\ii\tau
\]
is a nontrivial zero with
\[
        0<u\le0.01,
        \qquad
        \tau\ge T.
\]
By \cref{lem:reflected-zero-pair}, the reflected point
\[
        s_+=\frac12+u+\ii\tau
\]
is also a nontrivial zero.  Hence, by \cref{lem:detector-vanishing},
\[
        \Dvec(u,\tau)=(0,0).
\]

On the other hand, by the two-channel decomposition
\cref{prop:two-channel-decomposition},
\[
        \Dvec(u,\tau)
        =
        \frac{\ii}{\tau}\Bvec(u,\tau)
        +
        \Rvec(u,\tau).
\]
Therefore
\[
        \|\Dvec(u,\tau)\|
        \ge
        \frac{\|\Bvec(u,\tau)\|}{\tau}
        -
        \|\Rvec(u,\tau)\|.
\]
By \cref{prop:certified-inner-bracket-vector},
\[
        \|\Bvec(u,\tau)\|\ge \gev.
\]
By \cref{cor:certified-residual-vector-estimate},
\[
        \|\Rvec(u,\tau)\|
        \le
        \Cinner
        \frac{\log\tau}{\tau^{3/2}}.
\]
Thus
\[
        \|\Dvec(u,\tau)\|
        \ge
        \frac{\gev}{\tau}
        -
        \Cinner
        \frac{\log\tau}{\tau^{3/2}}.
\]
The right-hand side is positive whenever
\[
        \sqrt{\tau}
        >
        \frac{\Cinner}{\gev}
        \log\tau.
\]

It remains to pass from the height inequality at \(T\) to every
\(\tau\ge T\).  Since
\[
        t\mapsto \frac{\sqrt t}{\log t}
\]
is increasing for \(t>e^2\), the inequality
\[
        \sqrt T
        >
        \frac{\Cinner}{\gev}
        \log T
\]
implies
\[
        \sqrt{\tau}
        >
        \frac{\Cinner}{\gev}
        \log\tau
\]
for every
\[
        \tau\ge T,
        \qquad
        T\ge e^2.
\]
Therefore
\[
        \|\Dvec(u,\tau)\|>0.
\]
This contradicts
\[
        \Dvec(u,\tau)=(0,0).
\]
Hence no such reflected off-critical zero pair exists.
\end{proof}

\begin{remark}[Why this is genuinely an inner-strip theorem]
The theorem does not rely on a one-point lower bound for
\[
        |\Beff(\sigma,\tau)|.
\]
Instead it uses the reflected vector lower bound
\[
        \|\Bvec(u,\tau)\|\ge\gev.
\]
This is why the theorem remains stable as \(u\to0\).  The divided odd channel
is precisely what captures the transverse derivative information lost by the
one-channel outer-strip bracket.
\end{remark}

% ============================================================
% 9. Height bridge and public numerical target
% ============================================================

\section{Height bridge and public numerical target}
\label{sec:height-bridge}

The certified height condition is
\[
        \sqrt T
        >
        \frac{\Cinner}{\gev}
        \log T.
\]
The public bracket-vector certificate uses the conservative value
\[
        \gev=0.2.
\]
Therefore the sufficient height condition becomes
\[
        \sqrt T
        >
        5\Cinner\log T.
\]

For the shared height
\[
        T=10^{30},
\]
this condition is equivalent to
\[
        \Cinner
        <
        0.2\frac{\sqrt{10^{30}}}{\log(10^{30})}.
\]
Since
\[
        \sqrt{10^{30}}=10^{15},
\]
and
\[
        \log(10^{30})=30\log10,
\]
it is enough that
\[
        \Cinner
        \le
        2.895296546022\times10^{12}.
\]

\begin{proposition}[Public height criterion for \(10^{30}\)]
\label{prop:height-criterion-1e30}
Assume the public bracket-vector lower bound
\[
        \gev=0.2.
\]
If the final inner certificate gives
\[
        \Cinner
        \le
        2.895296546022\times10^{12},
\]
then
\[
        T=10^{30}
\]
is an admissible certified inner-strip height.
\end{proposition}

\begin{proof}
Substituting
\[
        T=10^{30}
\]
and
\[
        \gev=0.2
\]
into the height condition gives
\[
        10^{15}
        >
        \frac{\Cinner}{0.2}
        30\log10.
\]
This is equivalent to
\[
        \Cinner
        <
        0.2\frac{10^{15}}{30\log10}.
\]
The right-hand side is
\[
        2.895296546022\ldots\times10^{12}.
\]
Thus the displayed upper bound for \(\Cinner\) implies the height inequality.
\end{proof}

\begin{corollary}[Certified inner strip at the shared height]
\label{cor:inner-strip-shared-height}
Assume the hypotheses of
\cref{thm:conditional-certified-inner-strip}.  Assume further that the final
public certificate records
\[
        \gev=0.2
\]
and
\[
        \Cinner
        \le
        2.895296546022\times10^{12}.
\]
Then there are no nontrivial zeros in
\[
        0<\left|\Re s-\frac12\right|\le0.01,
        \qquad
        \Im s\ge10^{30}.
\]
\end{corollary}

\begin{proof}
This follows immediately from
\cref{thm:conditional-certified-inner-strip} and
\cref{prop:height-criterion-1e30}.
\end{proof}

\begin{remark}[If the final constant is larger]
If the final certificate gives
\[
        \Cinner>2.895296546022\times10^{12},
\]
then the inner-strip theorem is not invalid.  One simply cannot use the shared
height \(10^{30}\) with the conservative value \(\gev=0.2\).  The certified
height must then be recomputed from
\[
        \sqrt T
        >
        \frac{\Cinner}{0.2}
        \log T.
\]
\end{remark}

% ============================================================
% 10. Public certificate interface
% ============================================================

\section{Public certificate interface}
\label{sec:public-certificate-interface}

The final inner-strip theorem is a certified theorem only when the public JSON
certificate exists and has the required status.  The purpose of this section is
to make the certificate contract explicit.

The final public certificate is
\[
        \texttt{INNER\_STRIP\_CERTIFICATE\_SAFE.json}.
\]
It must have
\[
        \texttt{status = "inner\_full\_certificate"}.
\]
It must also record the following fields:
\[
        \texttt{u0},
        \qquad
        \texttt{g0\_bracket\_vector\_lower},
\]
\[
        \texttt{C\_even\_residual\_upper},
        \qquad
        \texttt{C\_sigma\_residual\_upper},
\]
\[
        \texttt{C\_inner\_vector\_upper},
        \qquad
        \texttt{height\_tested},
\]
and
\[
        \texttt{height\_condition\_holds}.
\]

The final certificate must depend on seven sector certificates:
\[
        {\rm band},
        \quad
        {\rm bd},
        \quad
        {\rm far},
        \quad
        {\rm end},
        \quad
        {\rm nonstat},
        \quad
        {\rm floor},
        \quad
        {\rm core}.
\]
Each sector certificate must have status
\[
        \texttt{proved}
\]
and must contain both
\[
        \texttt{C\_value\_upper\_safe}
\]
and
\[
        \texttt{C\_sigma\_derivative\_upper\_safe}.
\]

\begin{definition}[Acceptable inner certificate]
\label{def:acceptable-inner-certificate}
An inner-strip certificate is acceptable if and only if:

\begin{enumerate}[label=(\roman*)]
    \item the final file
    \[
            \texttt{INNER\_STRIP\_CERTIFICATE\_SAFE.json}
    \]
    has
    \[
            \texttt{status = "inner\_full\_certificate"};
    \]

    \item the bracket-vector lower field satisfies
    \[
            \texttt{g0\_bracket\_vector\_lower}\ge0.2;
    \]

    \item every sector dependency has status
    \[
            \texttt{proved};
    \]

    \item every sector dependency contains both an ordinary residual constant
    and a \(\sigma\)-derivative residual constant;

    \item the final file records
    \[
            \texttt{height\_condition\_holds = true}.
    \]
\end{enumerate}
\end{definition}

\begin{remark}[No theorem without the final status]
A collection of partial certificates does not certify the inner-strip theorem.
The theorem is certified only after the strict merger produces the final status
\[
        \texttt{inner\_full\_certificate}.
\]
This prevents accidental use of a certificate missing one of the derivative
sectors.
\end{remark}

\subsection{Reproduction commands}

The expected public workflow is the following.  First create the inner
certificate directory:
\begin{verbatim}
mkdir -p certs_inner
\end{verbatim}

Run the bracket-vector verifier:
\begin{verbatim}
python scripts/cert_inner_bracket_vector.py \
  > certs_inner/inner_bracket_vector.raw.json
\end{verbatim}

Then run the seven derivative-aware sector verifiers.  Their final safe outputs
must have the following form:
\begin{verbatim}
certs_inner/band_sector.safe.json
certs_inner/bd_sector.safe.json
certs_inner/far_sector.safe.json
certs_inner/end_sector.safe.json
certs_inner/nonstat_sector.safe.json
certs_inner/floor_sector.safe.json
certs_inner/core_sector.safe.json
\end{verbatim}

Finally run the strict merger:
\begin{verbatim}
python scripts/merge_inner_certificate_safe.py \
  --inputs 'certs_inner/*_sector.safe.json' \
  --height 1e30 \
  --g0 0.2 \
  --output INNER_STRIP_CERTIFICATE_SAFE.json
\end{verbatim}

The theorem at height \(10^{30}\) is certified only if the final output has
\[
        \texttt{status = "inner\_full\_certificate"}.
\]

% ============================================================
% 11. Combined consequence with the outer-strip theorem
% ============================================================


\section{Combined consequence with the certified outer strip}
\label{sec:combined-consequence}

The certified outer-strip theorem gives the zero-free region
\[
        0.01\le\Re s\le0.49,
        \qquad
        \Im s\ge10^{30},
\]
and, by reflection,
\[
        0.51\le\Re s\le0.99,
        \qquad
        \Im s\ge10^{30}.
\]
The inner-strip theorem certified in this paper covers
\[
        0<\left|\Re s-\frac12\right|\le0.01,
        \qquad
        \Im s\ge10^{30},
\]
provided the final inner certificate satisfies the public height criterion.

Thus, after both the outer-strip and inner-strip certificates are complete, the
combined certified region is
\[
        0.01\le\Re s\le0.99,
        \qquad
        \Re s\ne\frac12,
        \qquad
        \Im s\ge10^{30}.
\]

\begin{corollary}[Combined certified off-critical exclusion above the shared height]
\label{cor:combined-off-critical-exclusion}
Assume the certified outer-strip theorem and assume the acceptable inner
certificate of \cref{def:acceptable-inner-certificate} verifies the height
\[
        T=10^{30}.
\]
Then the Riemann zeta function has no nontrivial zeros with
\[
        0.01\le\Re s\le0.99,
        \qquad
        \Re s\ne\frac12,
        \qquad
        \Im s\ge10^{30}.
\]
\end{corollary}

\begin{proof}
The outer-strip theorem covers
\[
        0.01\le\Re s\le0.49
\]
and its reflected strip
\[
        0.51\le\Re s\le0.99.
\]
The inner-strip theorem covers
\[
        0<\left|\Re s-\frac12\right|\le0.01,
\]
which is precisely
\[
        0.49\le\Re s<\frac12
        \quad\text{and}\quad
        \frac12<\Re s\le0.51.
\]
Together these regions cover
\[
        0.01\le\Re s\le0.99,
        \qquad
        \Re s\ne\frac12.
\]
\end{proof}

\begin{remark}[Extreme strips near \(0\) and \(1\)]
The statement above does not claim a certified exclusion in
\[
        0<\Re s<0.01
        \quad\text{or}\quad
        0.99<\Re s<1.
\]
Those extreme strips would require either a separate certified outer-strip run
with a smaller left endpoint, or an independent zero-free argument covering
those regions.
\end{remark}

\begin{remark}[This is not a finite-height verification below \(10^{30}\)]
The combined theorem is a high-height off-critical exclusion.  It does not by
itself verify zeros below \(10^{30}\).  To turn the combined high-height
exclusion into a full verification statement up to the critical line, one would
also need an independent finite-height verification up to the same height.
No such finite-height verification is claimed in this paper.
\end{remark}

\begin{remark}[No claim about zeros on the critical line]
The theorem excludes off-critical zeros in the certified high-height region.
It does not exclude zeros on the critical line.  The critical line is allowed
and is not part of the contradiction.
\end{remark}

% ============================================================
% Bibliography
% ============================================================

\begin{thebibliography}{9}

\bibitem{OuterStructural}
V\'ictor Blanco Bataller.
\newblock \emph{A Mellin-Symmetric Quotient Obstruction and a Certified Outer-Strip Exclusion for the Riemann Zeta Function}.
\newblock Structural outer-strip paper, 2026.

\bibitem{OuterCertifier}
V\'ictor Blanco Bataller.
\newblock \emph{A Certified Outer-Strip Exclusion for the Zeta Quotient Obstruction}.
\newblock Companion certification paper, 2026.

\bibitem{CertRepo}
V\'ictor Blanco Bataller.
\newblock \emph{RH outer-strip certifier repository}.
\newblock Public verification repository, 2026.
\newblock \RepoURL.
\newblock Commit: \texttt{TO-BE-INSERTED}.

\bibitem{Titchmarsh}
E. C. Titchmarsh.
\newblock \emph{The Theory of the Riemann Zeta-Function}.
\newblock Second edition, revised by D. R. Heath-Brown.
\newblock Oxford University Press, 1986.

\bibitem{Ivic}
A. Ivi\'c.
\newblock \emph{The Riemann Zeta-Function}.
\newblock Dover Publications, 2003.

\bibitem{Edwards}
H. M. Edwards.
\newblock \emph{Riemann's Zeta Function}.
\newblock Dover Publications, 2001.

\bibitem{JohanssonArb}
Fredrik Johansson.
\newblock Arb: Efficient arbitrary-precision midpoint-radius interval arithmetic.
\newblock \emph{IEEE Transactions on Computers}, 66(8):1281--1292, 2017.
\newblock DOI: \texttt{10.1109/TC.2017.2690633}.

\bibitem{FLINT}
The FLINT team.
\newblock \emph{FLINT: Fast Library for Number Theory}.
\newblock \url{https://flintlib.org}.

\bibitem{PythonFlint}
The python-flint developers.
\newblock \emph{python-flint: Python bindings for FLINT and Arb}.
\newblock \url{https://github.com/flintlib/python-flint}.

\end{thebibliography}

\end{document}